\documentclass[12pt,a4paper]{article}

\usepackage[a4paper,margin=1in]{geometry}
\usepackage{titlesec}
\usepackage{enumitem}
\usepackage{longtable}
\usepackage{array}
\usepackage{hyperref}

\title{Distributed Execution Platform Architecture}
\author{Project Team}
\date{\today}

\begin{document}

\maketitle

\tableofcontents
\newpage

\section{Introduction}

This document describes the architecture and responsibilities of the distributed execution platform. 
The system is composed of three major entities:

\begin{itemize}
    \item Worker Agent
    \item Controller Agent
    \item Master Node
\end{itemize}

Each component is divided into multiple threads or services responsible for computation, monitoring, orchestration, communication, and fault tolerance.

\section{System Overview}

The platform is designed to distribute workloads across multiple machines in a network. 
The master node distributes tasks, worker agents execute them, and controller agents maintain cluster state and monitor failures.

The architecture supports:

\begin{itemize}
    \item Parallel workload execution
    \item Dynamic node monitoring
    \item Fault tolerance
    \item Heartbeat-based liveness detection
    \item Distributed orchestration
    \item Reliable network communication
\end{itemize}

\newpage

\section{Worker Agent}

The Worker Agent is responsible for executing workloads assigned by the master node and reporting execution results.

\subsection{Execution Thread}

The execution thread is responsible for handling and executing workloads received from the master.

\subsubsection{Responsibilities}

\begin{itemize}
    \item Receive workloads from the network thread.
    \item Create a dedicated thread for each incoming workload.
    \item Identify the exact code segment to execute from the complete program source file.
    \item Execute the extracted code using the provided input data.
    \item Capture execution output and status.
    \item Send execution results to the network thread for transmission.
\end{itemize}

\subsubsection{Execution Flow}

\begin{enumerate}
    \item Receive workload.
    \item Create execution thread.
    \item Parse source file.
    \item Locate executable section.
    \item Execute code.
    \item Collect results.
    \item Forward results to network process.
\end{enumerate}

\subsection{Monitoring Thread}

The monitoring thread is responsible for collecting system statistics and hardware information.

\subsubsection{Responsibilities}

\begin{itemize}
    \item Read CPU information using commands such as \texttt{lscpu}.
    \item Monitor memory usage.
    \item Monitor disk usage.
    \item Monitor network utilization.
    \item Collect system load metrics.
    \item Send monitoring data periodically to the network thread.
\end{itemize}

\subsubsection{Example Commands}

\begin{itemize}
    \item \texttt{lscpu}
    \item \texttt{free -m}
    \item \texttt{df -h}
    \item \texttt{top}
    \item \texttt{uptime}
\end{itemize}

\subsection{Network Thread}

The network thread handles all communication between the worker and the distributed network.

\subsubsection{Responsibilities}

\begin{itemize}
    \item Poll local message queues for outgoing messages.
    \item Send messages to remote nodes.
    \item Receive incoming messages.
    \item Dispatch incoming messages to the appropriate local service.
    \item Confirm message delivery.
    \item Retry transmission when failures occur.
    \item Maintain a queue of unsent messages.
\end{itemize}

\subsubsection{Reliability Features}

\begin{itemize}
    \item Message acknowledgments.
    \item Retry mechanisms.
    \item Persistent unsent-message queue.
    \item Timeout detection.
    \item Failure recovery support.
\end{itemize}

\newpage

\section{Controller Agent}

The Controller Agent maintains a local representation of the distributed system state.

\subsection{Purpose}

The controller agent acts as an intermediary between worker nodes and the master node by maintaining node status, monitoring failures, and handling coordination tasks.

\subsection{Local State Storage}

The controller stores:

\begin{itemize}
    \item Node identifiers
    \item CPU utilization
    \item Memory usage
    \item Availability status
    \item Heartbeat timestamps
    \item Active workload information
\end{itemize}

\subsection{Network Thread}

The network thread inside the controller handles all communication between controller services and the distributed cluster.

\subsection{State Receiver}

The state receiver is responsible for receiving and updating node information.

\subsubsection{Responsibilities}

\begin{itemize}
    \item Receive ``hello'' messages from newly joined nodes.
    \item Register new agents into local state storage.
    \item Receive heartbeat messages from agents.
    \item Update node status continuously.
    \item Detect inactive nodes through heartbeat timeout.
\end{itemize}

\subsection{Master Servant}

The Master Servant acts as a service provider to the master node.

\subsubsection{Responsibilities}

\begin{itemize}
    \item Listen for requests from the master.
    \item Provide local state information.
    \item Return status reports for all known agents.
    \item Send monitoring data when requested.
\end{itemize}

\subsection{Panne Detection Service}

The Panne service is responsible for failure detection and overload monitoring.

\subsubsection{Responsibilities}

\begin{itemize}
    \item Monitor local states for overloaded agents.
    \item Detect crashed or unreachable nodes.
    \item Inform the master node when failures occur.
    \item Detect master node failure.
    \item Trigger election procedures when the master is unavailable.
    \item Assist replacement master initialization.
\end{itemize}

\newpage

\section{Master Node}

The Master Node coordinates workload execution across the distributed system.

\subsection{Network Thread}

The network thread handles all communications between the master and distributed agents.

\subsubsection{Responsibilities}

\begin{itemize}
    \item Receive execution requests.
    \item Send workloads to workers.
    \item Receive execution results.
    \item Communicate with controller agents.
    \item Handle acknowledgments and retries.
\end{itemize}

\subsection{Execution Master}

The Execution Master coordinates distributed execution.

\subsubsection{Responsibilities}

\begin{itemize}
    \item Execute local master-side logic.
    \item Prepare workloads.
    \item Coordinate distributed execution.
    \item Interface with the parser and orchestrator.
\end{itemize}

\subsection{Parser}

The parser transforms annotated source code into executable distributed code.

\subsubsection{Responsibilities}

\begin{itemize}
    \item Parse annotated programs.
    \item Detect distributed execution annotations.
    \item Replace annotations with appropriate execution code.
    \item Generate workload segments.
    \item Prepare tasks for orchestration.
\end{itemize}

\subsection{Orchestrator}

The orchestrator distributes workloads among available worker agents.

\subsubsection{Responsibilities}

\begin{itemize}
    \item Receive tasks from the execution master.
    \item Query controller agents for available nodes.
    \item Determine workload distribution strategy.
    \item Split workloads proportionally.
    \item Dispatch workloads to workers.
    \item Monitor execution progress.
    \item Handle task redistribution on failures.
\end{itemize}

\newpage

\section{Fault Tolerance Mechanisms}

The system incorporates several fault tolerance techniques.

\subsection{Heartbeat Monitoring}

Nodes periodically send heartbeat messages to indicate availability.

\subsection{Retry Mechanisms}

All communications include acknowledgment checks and retry procedures.

\subsection{Leader Election}

If the master node fails:

\begin{itemize}
    \item Controller agents detect failure.
    \item Election procedure is triggered.
    \item New master node is selected.
    \item Cluster state is restored.
\end{itemize}

\subsection{Message Persistence}

Unsent messages are stored locally until successful delivery.

\section{Concurrency Model}

The platform uses a multi-threaded architecture.

\subsection{Thread Isolation}

Each major responsibility runs independently:

\begin{itemize}
    \item Execution Threads
    \item Monitoring Threads
    \item Network Threads
    \item Failure Detection Threads
    \item Orchestration Threads
\end{itemize}

\subsection{Advantages}

\begin{itemize}
    \item Parallel execution
    \item Improved responsiveness
    \item Better scalability
    \item Isolation of failures
\end{itemize}

\newpage

\section{Team Responsibilities}

\begin{longtable}{|p{5cm}|p{8cm}|}
\hline
\textbf{Member} & \textbf{Assigned Component} \\
\hline
Tchami & Monitoring + Agent Initialization \\
\hline
Farelle & State Receiver \\
\hline
Ngonga & Network Layer \\
\hline
Tchapda & Worker Execution Module \\
\hline
Nzieleu & Orchestration Module \\
\hline
Felix & Master Servant \\
\hline
Kouam & Panne Detection Service \\
\hline
Bala & Execution Master \\
\hline
Ulrich & Execution Master \\
\hline
Tagne & Panne Detection Service \\
\hline
Roy & Parser \\
\hline
\end{longtable}

\section{Conclusion}

This architecture provides a scalable and fault-tolerant distributed execution platform capable of:

\begin{itemize}
    \item Executing distributed workloads efficiently
    \item Monitoring cluster resources
    \item Detecting node failures
    \item Dynamically orchestrating execution
    \item Recovering from communication and node failures
\end{itemize}

The separation of responsibilities between Worker Agents, Controller Agents, and the Master Node ensures modularity, scalability, and maintainability.

\end{document}